% ============================================================
%  Equações e Inequações do 1º Grau — Apostila Premium | PratiqueJá
%  Engine: XeLaTeX
% ============================================================
\documentclass[12pt]{article}
\usepackage{../../pratiqueja}
\usepackage{tikz}

% \hlm — destaque de resultado em modo-math (cor sem bold)
\newcommand{\hlm}[1]{{\color{darkblue}#1}}

\setsubject{Equações e Inequações do 1º Grau}

\begin{document}
\pagenumbering{arabic}
\setstretch{1.2}
\color{bodytext}

\heroheader{Equações do 1º Grau}%
           {Resolução, Casos Especiais e Inequações}%
           {Matemática · Intermediário}

\setlength{\columnsep}{18pt}
\setlength{\columnseprule}{0.5pt}
\def\columnseprulecolor{\color{exborder}}

\begin{multicols}{2}

%─────────────────────────────────────────────────────────────
\TW{Def}{badgepurple}{Definição}{Igualdade com incógnita, grau e solução}

\regra{Uma \textbf{equação} é uma igualdade com \textbf{incógnitas}.
O valor que torna a igualdade verdadeira é a \textbf{raiz} (solução).
O \textbf{grau} é o maior expoente da incógnita.}

\small\color{bodytext}%
No 1º grau a incógnita tem expoente $1$. Forma geral:

\formula{$ax + b = 0 \quad (a \neq 0) \;\Rightarrow\; x = -\dfrac{b}{a}$}

\exemp{%
  $2x - 3 = 7$ \qquad $6x - 1 = 4$\\[2pt]
  $2 - 5x = 9$ \qquad $3x - x - 5 = 3$%
}

%─────────────────────────────────────────────────────────────
\TW{Met}{darkblue}{Como Resolver}{Transposição e roteiro}

\regra{Isole a incógnita: termos com $x$ de um lado, constantes do
outro. Ao \textbf{trocar de lado}, inverta o operador. Ao final,
\textbf{verifique} substituindo a raiz.}

\SH{Regras de transposição}
\exemp{%
  $+$ vira $-$ \qquad $-$ vira $+$\\[2pt]
  $\times$ vira $\div$ \qquad $\div$ vira $\times$%
}

\SH{Roteiro}
\obs{\textbf{1.} Expandir parênteses. \;\textbf{2.} Transpor: $x$ à
esquerda, constantes à direita. \;\textbf{3.} Somar semelhantes.
\;\textbf{4.} Dividir pelo coeficiente de $x$. \;\textbf{5.}
Verificar.}

%─────────────────────────────────────────────────────────────
\TW{Ex}{badgegreen}{Exemplos Básicos}{Incógnita em um lado, com verificação}

\exemp{%
  $4x - 6 = 2$\\
  $4x = 2 + 6 = 8$ \;{\footnotesize($-6$ vira $+6$)}\\
  $x = \dfrac{8}{4} = \hlm{2}$\\
  {\footnotesize Verif.: $4\cdot2 - 6 = 2$ \ok}%
}

\exemp{%
  $3x - x - 5 = 3$\\
  $2x = 3 + 5 = 8$\\
  $x = \hlm{4}$\\
  {\footnotesize Verif.: $12 - 4 - 5 = 3$ \ok}%
}

%─────────────────────────────────────────────────────────────
\TW{2x}{darkblue}{Incógnita nos Dois Lados}{Juntar os $x$ do mesmo lado}

\regra{Transporte todos os termos com $x$ para a esquerda e as
constantes para a direita, depois simplifique.}

\exemp{%
  $3x + 2 = x + 8$\\
  $3x - x = 8 - 2 \;\Rightarrow\; 2x = 6$\\
  $x = \hlm{3}$ \quad{\footnotesize($11 = 11$ \ok)}%
}

\exemp{%
  $5x - 4 = 2x + 11$\\
  $5x - 2x = 11 + 4 \;\Rightarrow\; 3x = 15$\\
  $x = \hlm{5}$ \quad{\footnotesize($21 = 21$ \ok)}%
}

%─────────────────────────────────────────────────────────────
\TW{$(\,)$}{darkblue}{Equações com Parênteses}{Expandir antes de transpor}

\regra{Com parênteses, \textbf{expanda primeiro} (distributiva); só
então transponha.}

\exemp{%
  $2(x + 3) = 10$\\
  $2x + 6 = 10 \;\Rightarrow\; 2x = 4$\\
  $x = \hlm{2}$ \quad{\footnotesize($2\cdot5 = 10$ \ok)}%
}

\exemp{%
  $3(2x - 1) = 2(x + 4)$\\
  $6x - 3 = 2x + 8$\\
  $4x = 11 \;\Rightarrow\; x = \hlm{\dfrac{11}{4}}$%
}

%─────────────────────────────────────────────────────────────
\TW{Frac}{badgepurple}{Equações com Frações}{Eliminar denominadores pelo MMC}

\regra{Multiplique \textbf{todos os termos} pelo MMC dos
denominadores para eliminar as frações.}

\exemp{%
  $\dfrac{x}{2} + 1 = 3$ \;{\footnotesize(MMC $=2$)}\\
  $x + 2 = 6 \;\Rightarrow\; x = \hlm{4}$\\
  {\footnotesize Verif.: $\frac{4}{2}+1 = 3$ \ok}%
}

\exemp{%
  $\dfrac{2x - 1}{3} = 5$ \;{\footnotesize(MMC $=3$)}\\
  $2x - 1 = 15 \;\Rightarrow\; 2x = 16$\\
  $x = \hlm{8}$%
}

%─────────────────────────────────────────────────────────────
\TW{!}{vered}{Casos Especiais}{Sem solução e infinitas soluções}

\regra{Se a incógnita \textbf{desaparece}, analise a igualdade: se
for \textbf{falsa}, não há solução; se \textbf{verdadeira}, há
infinitas.}

\exemp{%
  \textbf{Impossível:}\\
  $2x + 3 = 2x + 7 \;\Rightarrow\; 0 = 4$ \nok\\
  Não tem solução: $\hlm{S = \emptyset}$\\[3pt]
  \textbf{Identidade:}\\
  $3(x+1) = 3x + 3 \;\Rightarrow\; 0 = 0$ \ok\\
  Infinitas soluções: $\hlm{S = \mathbb{R}}$%
}

%─────────────────────────────────────────────────────────────
\TW{Prob}{badgegreen}{Problema Contextualizado}{Montar e resolver a partir do enunciado}

\regra{\textbf{Estratégia:} (1) defina a incógnita; (2) traduza o
enunciado; (3) resolva; (4) interprete.}

\exemp{%
  \textbf{P1:} a soma de dois consecutivos é $15$.\\
  $x + (x+1) = 15 \;\Rightarrow\; 2x + 1 = 15$\\
  $2x = 14 \;\Rightarrow\; x = 7$\\
  $\hlm{\text{os números são 7 e 8}}$%
}

\exemp{%
  \textbf{P2:} um número $+5$ é o dobro dele $-3$.\\
  $x + 5 = 2x - 3$\\
  $5 + 3 = 2x - x \;\Rightarrow\; x = \hlm{8}$%
}

%─────────────────────────────────────────────────────────────
\TW{Ineq}{badgepurple}{Inequações do 1º Grau}{Desigualdades e intervalo solução}

\regra{Uma \textbf{inequação} é uma desigualdade com incógnita. Sua
solução é um \textbf{intervalo} na reta real, não um ponto.}

\SH{Símbolos}
\obs{$<$ menor \quad $>$ maior \quad $\leq$ menor/igual \quad
$\geq$ maior/igual}

\SH{Regra de inversão}
\regra{Ao multiplicar ou dividir por um número \textbf{negativo}, o
sentido da desigualdade \textbf{inverte}: $a < b \Rightarrow -a > -b$.
Somar/subtrair nunca muda o sentido.}

\exemp{%
  $3x - 2 > 7 \;\Rightarrow\; 3x > 9$\\
  $\hlm{x > 3} \;\Rightarrow\; x \in (3,\,+\infty)$\\[3pt]
  $-2x + 1 \leq 5 \;\Rightarrow\; -2x \leq 4$\\
  $\div(-2)$ inverte: $\hlm{x \geq -2} \;\Rightarrow\; x \in [-2,\,+\infty)$\\[3pt]
  $5 - 3x \leq 2x + 10 \;\Rightarrow\; -5x \leq 5$\\
  $\div(-5)$ inverte: $\hlm{x \geq -1}$%
}

\SH{Notação de intervalo}
\begin{center}\vspace{1pt}
\begin{tikzpicture}[scale=1]
% reta superior: fechado [-2, +inf)
\draw[line width=1pt, gray!50] (-2.3,0.45) -- (2.5,0.45);
\foreach \x in {-2,-1,0,1,2}{\draw[gray!50] (\x,0.38)--(\x,0.52);}
\draw[line width=2pt, darkblue] (-1,0.45) -- (2.5,0.45);
\fill[darkblue] (-1,0.45) circle (2.4pt);
\node[font=\scriptsize, below] at (-1,0.32){$-2$};
\node[font=\scriptsize, darkblue] at (1.9,0.78){$[-2,+\infty)$};
% reta inferior: aberto (3, +inf) -> aqui ilustro (a) aberto
\draw[line width=1pt, gray!50] (-2.3,-0.55) -- (2.5,-0.55);
\foreach \x in {-2,-1,0,1,2}{\draw[gray!50] (\x,-0.62)--(\x,-0.48);}
\draw[line width=2pt, badgegreen!60!black] (0,-0.55) -- (2.5,-0.55);
\draw[badgegreen!60!black, line width=1.5pt, fill=white] (0,-0.55) circle (2.4pt);
\node[font=\scriptsize, below] at (0,-0.68){$a$};
\node[font=\scriptsize, badgegreen!50!black] at (1.7,-0.25){$(a,+\infty)$};
\end{tikzpicture}
\end{center}

\obs{Bolinha \textbf{cheia} / colchete $[\,]$ = extremo \emph{incluído}
$(\leq,\geq)$. Bolinha \textbf{vazia} / parêntese $(\,)$ = extremo
\emph{excluído} $(<,>)$.}

\end{multicols}

\end{document}
